\section{Academic landscape and evidence-qualified directory}
\label{sec:directory}
The directory distinguishes demonstrated direct hot-glass deposition from adjacent glass research. Institutional inclusion follows the requested search scope and should not be read as a ranking. A historical project establishes activity at its publication date, not current equipment availability or personnel. ``Not verified'' records a limit of this investigation; it is not proof that a group has never worked on the topic.

\begin{longtable}{P{0.18\textwidth}P{0.22\textwidth}P{0.30\textwidth}P{0.20\textwidth}}
\caption{Global academic research directory. Named people are documented researchers or contacts, not necessarily current laboratory directors. Bibliography entries provide primary DOI/URL links.}\label{tab:academic_directory}\\
\toprule
Institution / laboratory & Researcher / apparatus & Verified contribution and route & Primary reference \\
\midrule
\endfirsthead
\toprule
Institution / laboratory & Researcher / apparatus & Verified contribution and route & Primary reference \\
\midrule
\endhead
MIT Mediated Matter and Glass Lab & Neri Oxman, John Klein, Peter Houk and collaborators; G3DP, G3DP2 & Direct reservoir extrusion; integrated heating and motion; architectural demonstrators. Historical project attribution. & \cite{acad_klein2015,acad_inamura2018} \\
TU Delft structural glass research & Telesilla Bristogianni, Faidra Oikonomopoulou; molten printer name not verified & Cast/recycled structural glass and fabrication research. Casting evidence must not be relabeled as direct AM. & Institutional publication directory \cite{acad_delft} \\
LLNL glass AM / optics & Rebecca Dylla-Spears, Du Nguyen, Joel Destino and collaborators; DIW research systems & Silica and mixed-oxide inks, subsequent thermal densification, optical characterization. Indirect route. & \cite{acad_nguyen2017,acad_destino2018} \\
Fraunhofer IKTS & Glass/ceramic processing teams; specific PI and printer name not verified & Glass-filled filament shaping followed by sintering; photomultiplier-related components. & Institute account \cite{acad_ikts} \\
Fraunhofer ISC with IOF & Project team; specific molten printer not verified & ORMOCER hybrid optical components; inorganic--organic polymer material must be distinguished from glass melt. & Institute account \cite{acad_isc} \\
Fraunhofer IPT & 3DHeat team; metal LPBF and wire deposition tooling & Additively manufactured metal molds with thermal control for glass molding. The tool is printed, not the glass. & 2025--2027 project \cite{acad_ipt} \\
KTH / Nobula & Chunxin Liu, Taras Oriekhov, Michael Fokine; laser filament research printer & Direct fused-silica filament deposition with symmetric heating and measured interlayer bonding. & \cite{acad_liu2022} \\
RISE Acreo with KTH / NTNU & Individual PI/apparatus assignment not verified & Institutional evidence of collaboration in laser processing and glass AM; no separate RISE printer specification inferred. & KTH project record \cite{acad_kthrise} \\
Universit\'e Laval, COPL & Yannick Ledemi, Youn\`es Messaddeq and collaborators; modified commercial extrusion printer & Direct arsenic-sulfide chalcogenide glass deposition. Separate LRIO work addresses nanocomposite photopolymerization. & \cite{acad_baudet2019,acad_lavalproject} \\
McGill University & Inorganic molten-glass PI/apparatus not verified & Search retrieved carbohydrate-glass bioprinting evidence, which is outside this review's inorganic-glass scope. & Repository distinction \cite{acad_mcgill} \\
Nagaoka University of Technology & Historical Takayuki Komatsu glass/crystal research; direct printer not verified & Glass/crystal and laser-processing research provides adjacent expertise; direct molten deposition claim remains unverified. & Institutional record \cite{acad_nagaoka} \\
Laser Zentrum Hannover / Quarzglas-Technologie Heinrich & Philipp von Witzendorff, Leonhard Pohl and collaborators; CO$_2$ deposition head & Direct quartz-fiber deposition using laser heating. & \cite{acad_witzendorff2018} \\
TU Darmstadt / Deakin & Matthias Seel, Nick Birbilis and collaborators; extruded soda-lime specimens & Toolpath-dependent mechanical behavior and four-point bend tests; useful 2026 structural evidence. & \cite{acad_chhadeh2026} \\
\bottomrule
\end{longtable}

\subsection{How to interpret the research network}
The most useful collaboration depends on the missing capability. Reservoir extrusion needs melt conditioning, refractory compatibility, motion coordination, and annealing capacity. Fused-silica filament deposition depends more directly on laser delivery, feedstock uniformity, and local bonding. Optical-quality ink routes instead place substantial demands on colloid formulation, drying, organic removal, densification, and finishing. These categories overlap in characterization but do not share interchangeable processing temperatures or material supply chains.

Two attribution errors deserve explicit correction. First, publication in proceedings hosted at TU Delft does not mean that the printer belongs to TU Delft: the 2024 recycled-glass extrusion study identifies Evenline and MIT authors.\cite{acad_stern2024} Second, the KTH article's affiliations identify KTH and Nobula; its editor's RISE affiliation is not evidence that RISE operated that apparatus. The separate institutional collaboration record establishes RISE involvement at project level, with a narrower evidentiary claim.\cite{acad_liu2022,acad_kthrise}

Laval is a verified direct-glass-extrusion contributor, but its arsenic-sulfide composition belongs to the chalcogenide family and uses a substantially different processing window from silicates.\cite{acad_baudet2019} Its later nanocomposite-resin project is an additional route, not evidence that the chalcogenide platform prints silica.\cite{acad_lavalproject} McGill and Nagaoka remain explicit search gaps for the requested direct molten inorganic-glass directory. Resolving them would require a primary apparatus paper, thesis, patent, or laboratory record identifying both the material and the deposition process; general expertise in glass or additive manufacturing is insufficient.

This evidence structure favors reproducible comparisons and appropriate research partnerships. A prospective collaboration brief should specify composition, size range, required optical function, and final thermal treatment before naming a preferred institution. Such a brief is more informative than a broad list of organizations associated with the phrase ``glass 3D printing.''
